\section{The three-way splitting in Hastings' approximate ground-state projector}
\label{sec:hastings_three_way_splitting}

% Hastings, arXiv:0705.2024v4, law16.tex, lines 763--790 (step one of the
% proof of Lemma 1); setting at lines 82--91, statement at lines 246--267.

The proof of Lemma~1 of \cite{Hastings2007AreaLaw} begins by splitting the
Hamiltonian at a cut into a left part, a buffer part and a right part, and by
centering each part in the ground vector. This section carries out that step
for a finite-range Hamiltonian of range $R$ on the chain of $L$ sites. The
source works with nearest-neighbour terms and a buffer at distance $l/3$ from
the cut; here the buffer width is an arbitrary natural number $m$. Norms are
operator norms.

\begin{definition}[Three-way splitting at a cut]
    \label{def:hastings_three_way_splitting}
    \lean{
        TNLean.MPS.AreaLaw.Hastings.leftSupports,
        TNLean.MPS.AreaLaw.Hastings.bufferSupports,
        TNLean.MPS.AreaLaw.Hastings.rightSupports,
        TNLean.MPS.AreaLaw.Hastings.sum_supports_eq_left_add_buffer_add_right,
        TNLean.MPS.AreaLaw.ChainHamiltonian.termSum,
        TNLean.MPS.AreaLaw.ChainHamiltonian.termSum_isHermitian,
        TNLean.MPS.AreaLaw.ChainHamiltonian.leftPart,
        TNLean.MPS.AreaLaw.ChainHamiltonian.bufferPart,
        TNLean.MPS.AreaLaw.ChainHamiltonian.rightPart,
        TNLean.MPS.AreaLaw.Hastings.groundCentered,
        TNLean.MPS.AreaLaw.Hastings.groundCentered_add,
        Matrix.IsHermitian.groundCentered,
        Matrix.IsHermitian.ofReal_re_inner
    }
    \leanok
    \uses{def:chain_model}
    Let $H=\sum_Xh_X$ be a finite-range Hamiltonian of range $R$ on the chain
    of $L$ sites, let $j$ be a cut, so that the sites $i<j$ lie on its left,
    and let $m$ be a natural number. The \emph{left part} $H_L$ is the sum of
    the terms $h_X$ with $X\subseteq[0,j-m)$, the \emph{right part} $H_R$ is the
    sum of the terms with $X\subseteq[j+m,L)$, and the \emph{buffer part}
    $H_B$ is the sum of all other terms. Each part is Hermitian. For a vector
    $\Psi_0$ and a matrix $M$, the \emph{centering} of $M$ in $\Psi_0$ is
    $\widetilde M=M-\operatorname{Re}\langle\Psi_0,M\Psi_0\rangle\,1$; it is
    additive in $M$, and Hermitian when $M$ is. For a Hermitian $M$ the
    expectation $\langle\Psi_0,M\Psi_0\rangle$ is real. These are the operators
    of \cite[lines 774--782]{Hastings2007AreaLaw}, with $l/3$ replaced by $m$.
\end{definition}

\begin{theorem}[The splitting is exact]
    \label{thm:hastings_three_way_splitting_sum}
    \lean{
        TNLean.MPS.AreaLaw.ChainHamiltonian.operator_eq_leftPart_add_bufferPart_add_rightPart,
        TNLean.MPS.AreaLaw.ChainHamiltonian.groundCentered_parts_sum,
        TNLean.MPS.AreaLaw.ChainHamiltonian.groundCentered_parts_sum_of_eigen,
        TNLean.MPS.AreaLaw.ChainHamiltonian.isGappedGroundState_groundCentered_parts_sum,
        TNLean.MPS.AreaLaw.Hastings.inner_groundCentered,
        TNLean.MPS.AreaLaw.Hastings.groundCentered_of_eigen
    }
    \leanok
    \uses{def:hastings_three_way_splitting}
    For every cut $j$ and every $m$, $H=H_L+H_B+H_R$. Let $\Psi_0$ be a unit
    vector with $H\Psi_0=E_0\Psi_0$. Then each centered part has expectation
    zero in $\Psi_0$, and
    \begin{align*}
        \widetilde H_L+\widetilde H_B+\widetilde H_R=H-E_0.
    \end{align*}
    If moreover $H-E_0\ge\Delta(1-|\Psi_0\rangle\langle\Psi_0|)$, then
    $\Psi_0$ is a gapped ground vector of
    $\widetilde H_L+\widetilde H_B+\widetilde H_R$ with energy zero and the
    same gap $\Delta$. With the convention of
    \cite[lines 766--767]{Hastings2007AreaLaw} that the ground energy is zero,
    this is the identity $H=H_L+H_B+H_R$ for the centered parts of
    \cite[lines 779--782]{Hastings2007AreaLaw}.
\end{theorem}
\begin{proof}\leanok
    Every support is left, right or buffer, and no support is both left and
    right, since a support is nonempty and $j-m\le j+m$. Summing the terms
    over the three classes gives $H=H_L+H_B+H_R$. Centering is additive, so
    the centered parts sum to $H-\operatorname{Re}\langle\Psi_0,H\Psi_0\rangle$,
    and $\langle\Psi_0,H\Psi_0\rangle=E_0$. The expectation of
    $\widetilde M$ is $\langle\Psi_0,M\Psi_0\rangle-
        \operatorname{Re}\langle\Psi_0,M\Psi_0\rangle=0$ for a Hermitian $M$.
    The gap condition is unchanged, since the shifted operator differs from
    $H$ by the scalar $E_0$.
\end{proof}

\begin{theorem}[Supports of the three parts]
    \label{thm:hastings_three_way_splitting_supports}
    \lean{
        TNLean.MPS.AreaLaw.ChainHamiltonian.isSupportedOn_termSum,
        TNLean.MPS.AreaLaw.ChainHamiltonian.isSupportedOn_leftPart,
        TNLean.MPS.AreaLaw.ChainHamiltonian.isSupportedOn_rightPart,
        TNLean.MPS.AreaLaw.ChainHamiltonian.isSupportedOn_bufferPart,
        TNLean.MPS.AreaLaw.Hastings.subset_window_of_mem_bufferSupports,
        Entropy.IsSupportedOn.groundCentered
    }
    \leanok
    \uses{def:hastings_three_way_splitting}
    The left part acts as the identity off $[0,j-m)$, the right part off
    $[j+m,L)$, and the buffer part off the window $[j-m-R,j+m+R)$. Centering
    preserves each of these supports. In particular the left and right parts
    lie at distance at least $m$ from the cut, for every $m$, as the parts
    $H_L$ and $H_R$ of \cite[lines 774--778]{Hastings2007AreaLaw} lie at
    distance $l/3$.
\end{theorem}
\begin{proof}\leanok
    A sum of terms supported in a set is supported in that set, and the
    identity is supported everywhere. A buffer support $X$ contains a site
    $x\ge j-m$ and a site $y<j+m$; since $X$ has diameter at most $R$, every
    site $z\in X$ satisfies $j-m-R\le z<j+m+R$.
\end{proof}

\begin{theorem}[Norm of the buffer part]
    \label{thm:hastings_three_way_splitting_norm}
    \lean{
        TNLean.MPS.AreaLaw.ChainHamiltonian.norm_termSum_le,
        TNLean.MPS.AreaLaw.Hastings.card_supports_containing_le,
        TNLean.MPS.AreaLaw.Hastings.card_bufferSupports_le,
        TNLean.MPS.AreaLaw.ChainHamiltonian.norm_bufferPart_le,
        TNLean.MPS.AreaLaw.Hastings.norm_groundCentered_le,
        TNLean.MPS.AreaLaw.ChainHamiltonian.norm_groundCentered_bufferPart_le
    }
    \leanok
    \uses{def:hastings_three_way_splitting,
        thm:hastings_three_way_splitting_supports}
    A sum of $k$ terms has norm at most $kJ$. At most $2^{2R+1}$ supports of
    range $R$ contain a given site, and there are at most
    $2(m+R)\,2^{2R+1}$ buffer supports. Hence, for $J\ge0$,
    \begin{align*}
        \|H_B\|\le2(m+R)\,2^{2R+1}J,
    \end{align*}
    linearly in the buffer width. For a unit vector $\Psi_0$, centering at
    most doubles the norm, so
    $\|\widetilde H_B\|\le4(m+R)\,2^{2R+1}J$.
\end{theorem}
\begin{proof}\leanok
    The first bound is the triangle inequality. A support containing a site
    $x$ lies in the $2R+1$ sites within distance $R$ of $x$, which have at most
    $2^{2R+1}$ subsets. Every buffer support contains a site of the window
    $[j-m-R,j+m+R)$, which has at most $2(m+R)$ sites. Finally
    $|\operatorname{Re}\langle\Psi_0,M\Psi_0\rangle|\le\|M\|$ for a unit
    vector $\Psi_0$, and the identity has norm at most one.
\end{proof}
